%%%%%%%%%%%%%%%%%
% This is an example CV created using altacv.cls (v1.1.5, 1 December 2018) written by
% LianTze Lim (liantze@gmail.com), based on the
% Cv created by BusinessInsider at http://www.businessinsider.my/a-sample-resume-for-marissa-mayer-2016-7/?r=US&IR=T
%
%% It may be distributed and/or modified under the
%% conditions of the LaTeX Project Public License, either version 1.3
%% of this license or (at your option) any later version.
%% The latest version of this license is in
%%    http://www.latex-project.org/lppl.txt
%% and version 1.3 or later is part of all distributions of LaTeX
%% version 2003/12/01 or later.
%%%%%%%%%%%%%%%%

%% If you are using \orcid or academicons
%% icons, make sure you have the academicons
%% option here, and compile with XeLaTeX
%% or LuaLaTeX.
% \documentclass[10pt,a4paper,academicons]{altacv}

%% Use the "normalphoto" option if you want a normal photo instead of cropped to a circle
% \documentclass[10pt,a4paper,normalphoto]{altacv}

\documentclass[11pt,a4paper,ragged2e,academicons]{altacv}
\usepackage{multicol}
% \usepackage{hyperref}
%% AltaCV uses the fontawesome and academicon fonts
%% and packages.
%% See texdoc.net/pkg/fontawecome and http://texdoc.net/pkg/academicons for full list of symbols. You MUST compile with XeLaTeX or LuaLaTeX if you want to use academicons.

% Change the page layout if you need to
\geometry{left=2cm,right=10cm,marginparwidth=6.8cm,marginparsep=1.2cm,top=1.25cm,bottom=1.25cm}

% Change the font if you want to, depending on whether
% you're using pdflatex or xelatex/lualatex
\ifxetexorluatex
  % If using xelatex or lualatex:
  \setmainfont{Carlito}
\else
  % If using pdflatex:
  \usepackage[utf8]{inputenc}
  \usepackage[T1]{fontenc}
  %\usepackage[default]{lato}
\fi
\usepackage{biblatex}
% Change the colours if you want to
\definecolor{VividPurple}{HTML}{000000}
\definecolor{SlateGrey}{HTML}{2E2E2E}
\definecolor{LightGrey}{HTML}{2E2E2E}
\colorlet{heading}{VividPurple}
\colorlet{accent}{VividPurple}
\colorlet{emphasis}{SlateGrey}
\colorlet{body}{LightGrey}

\definecolor{linkcol}{HTML}{0A58CA}

\usepackage{hyperref}

% Change the bullets for itemize and rating marker
% for \cvskill if you want to
\renewcommand{\itemmarker}{{\small\textbullet}}
\renewcommand{\ratingmarker}{\faCircle}

%% sample.bib contains your publications
\addbibresource{sample.bib}

\hypersetup{
    colorlinks,
    allcolors=linkcol
}

\preto\fullcite{\AtNextCite{\defcounter{maxnames}{99}}}


\begin{document}
\name{Nicole Sandra-Yaffa Dumont}
\tagline{}
% Cropped to square from https://en.wikipedia.org/wiki/Marissa_Mayer#/media/File:Marissa_Mayer_May_2014_(cropped).jpg, CC-BY 2.0
%\photo{3.3cm}{profile.jpg}
\personalinfo{%
  % Not all of these are required!
  % You can add your own with \printinfo{symbol}{detail}
   \location{Z\"urich, Switzerland}
   \email{Email:  \href{mailto:nidumon@ini.uzh.ch}{nidumon@ini.uzh.ch}}
\homepage{Website: \href{http://nsdumont.github.io/}{nsdumont.github.io}} \\
  \linkedin{LinkedIn: \href{https://www.linkedin.com/in/nicole-dumont/}{nicole-dumont}} 
   \github{Github: \href{https://github.com/nsdumont}{nsdumont} }
   %\homepage{compneuro.uwaterloo.ca/people/nicole-dumont.html} % .
   \scholar{Scholar: \href{https://scholar.google.com/citations?user=2hnh9gkAAAAJ&hl=en}{2hnh9gkAAAAJ}}
   \orcid{ORCID: \href{https://orcid.org/0000-0001-5041-1527}{ 0000-0001-5041-1527}} % 
}

%% Make the header extend all the way to the right, if you want.
\begin{fullwidth}
\makecvheader
% \end{fullwidth}

%% Depending on your tastes, you may want to make fonts of itemize environments slightly smaller
%\AtBeginEnvironment{itemize}{\small}

%% Provide the file name containing the sidebar contents as an optional parameter to \cvsection.
%% You can always just use \marginpar{...} if you do
%% not need to align the top of the contents to any
%% \cvsection title in the "main" bar.


\cvsection{Academic Appointments}

\cvevent{Postdoctoral Researcher}{Institute of Neuroinformatics (University of Zurich \& ETH Zurich)}{Dec 2025 -- Present}{Z\"urich, Switzerland}
\begin{itemize}
\item \textbf{Supervisor:} Prof. Katharina Wilmes
\item Awarded the \textbf{UZH Postdoc Grant}, beginning December 2026.
\item  \textbf{Predictive Coding (PC) \& Models of Perception and Learning in Cortex }
\begin{itemize}
    \item Comparing prediction error modulation by learned uncertainty signals and how this mechanism differs from top-down, goal-conditioned attentional modulation. 
    \item Developing PC networks for multi-sensory integration to better understand integration circuits in the brain. 
\end{itemize}
\end{itemize}

\cvsection{Education}

\cvevent{Ph.D in Computer Science \\ \normalsize Diploma in Theoretical Neuroscience}{University of Waterloo}{Sept 2019 -- April 2025}{Waterloo, Canada}
\begin{itemize}
\item \textbf{Supervisors:} Prof. Chris Eliasmith and Prof. Jeff Orchard
\item \textbf{Dissertation:} \emph{Symbols, Dynamics, and Maps: A Neurosymbolic Approach to Spatial Cognition} \href{https://hdl.handle.net/10012/21501}{[Link]}
    \item \textbf{Neuro-Symbolic Modeling for Spatial Representation \& Mapping}
\begin{itemize}\setlength\itemsep{0em}
        \item Developed Spatial Semantic Pointers (SSPs) as a grid cell-inspired representation of continuous features, enabling the integration of symbolic structure and uncertainty in biologically plausible neural models.
        \item Built spiking neural network models of path integration, cognitive mapping, and SLAM that use neurosymbolic methods to bind continuous and discrete features to create compositional embeddings. % leveraging SSPs for robust spatial inference on neuromorphic hardware.
     \end{itemize}
\item \textbf{Reinforcement Learning (RL) \& Action Selection}
\begin{itemize}\setlength\itemsep{0em}
\item Designed biologically plausible online learning rules for RL in continuous-time settings using State Space Model-based memory systems.
% \item Improved sample efficiency during navigation using SSPs and structured scene representations as inputs to deep RL frameworks.
\item Incorporated SSPs into Bayesian Optimization and deep RL frameworks, modeling predictive inference and decision-making under uncertainty.

%The Bayesian Optimization work was an INRC-funded project.

% in complex, multi-agent environments.
%     \item Achieved substantial gains in sample efficiency and compute time by leveraging structured internal models, mirroring cortical strategies for active exploration.
% \en
% \item Integrated SSPs into Bayesian Optimization frameworks for active exploration, achieving up to 175x reduction in compute time in multi-agent environments.
    \end{itemize}
    % \item \textbf{Interdisciplinary:} Research at the crossroads of cognitive science and computational neuroscience that advanced models integrating symbolic structure with neural dynamics to explain spatial cognition.
 \end{itemize}
 \vspace{-2pt}
\divider

\smallskip



\cvevent{Computational Mathematics (Masters of Mathematics -- Co-op Program)}{University of Waterloo}{Sept 2017 -- April 2019}{Waterloo, Canada}
% \begin{itemize}
% \item \textbf{Supervisor:} Prof. Thomas F. Coleman 
% \item \textbf{Thesis Project:} \emph{The Optimal Social Cost of Carbon \& The Cost of Inaction}
    % \item Developed and analyzed a robust optimization algorithm for carbon pricing models, integrating applied mathematics methods with real-world policy modeling.
% \item Applied automatic differentiation and Quasi-Newton methods to improve scalability, parameter sensitivity analysis, and robust optimization under uncertainty.
    %automatic differentiation
% and Quasi-Newton optimization to the EZ model, we significantly improve its computational
% performance. This improvement allows for increased scalability, parameter sensitivity anal-
% ysis, robust optimization, and evaluations of suboptimal policies
    % \item Cumulative average of 91 \% in course work, which included courses on optimization, computational statistics, numerical analysis, PDEs, and computational neuroscience.
% \end{itemize}
% \begin{itemize}
% \item Cumulative average of 91 \% 
% \item Completed a masters research paper on robust optimization of an asset pricing model used to price carbon emissions.
% \item Courses on optimization, computational statistics, numerical analysis, PDEs, and computational neuroscience.
% \end{itemize}


\divider
\smallskip

\cvevent{Honours Mathematics and Physics (Bachelors of Science) }{McMaster University}{Sept 2012 -- April 2017}{Hamilton, Canada}
%\begin{itemize}
 %\item Graduated with distinction. Courses on statistical mechanics, dynamical systems, mathematical biology, quantum computing, and scientific computation.
% \item Cumulative average of 10.5/12.0 (GPA 3.8)
%\end{itemize}

%\divider
 % \vspace{-10pt}
\cvsection{Work Experience}
% \smallskip


\cvevent{Research Associate}{Computational Neuroscience Research Group, University of Waterloo}{April 2025 -- Nov 2025}{Waterloo, Canada}
\begin{itemize}
    \item Contributed to the incorporation of spatial reasoning and embodiment in Spaun 3.0, a large-scale functional brain model.
\end{itemize}

%\pagebreak
\divider

\cvevent{Neuromorphic Scientist}{Applied Brain Research}{Jan 2023 -- Jan 2024}{Waterloo, Canada}
\begin{itemize}
    \item Integrated a biologically-inspired SLAM algorithm into a MuJoCo simulation environment and deployed it on low-power neuromorphic hardware (Intel Loihi chip).
\end{itemize}
% 

\divider
% \smallskip

\cvevent{Research Associate}{Cayuga Research}{May 2018 -- Aug 2019}{Waterloo, Canada}
 \begin{itemize}
     \item Consulting work focused on developing optimization methods and data-driven solutions for industrial applications.
%     % \item Built a global optimization toolbox in MATLAB.
% %\item Created prototype flight path optimization software achieving ~5\% fuel savings compared to real flights.
% %\item Designed and tested chiller plant optimization models, producing ~4\% energy savings.
 \end{itemize}
%     \item Projects included development of a global optimization toolbox in Matlab, prototype flight path optimization software, and chiller plant optimization \& modelling.
% \end{itemize}
% \begin{itemize}
% \item Worked as a part of a team for consulting work focused on the development and implementation of advanced optimization methods and data driven solutions to industrial problems. 
% %\item Developed a global optimization toolbox in Matlab
% \item Built prototype flight path optimization software able to plan flights that save up to 5\% in fuel costs compared to real commercial flights. 
% \item Worked on a chiller plant optimization problem, developing data-driven models of the plant components and an optimization method that produced 4\% energy savings.
% \end{itemize}
% \vspace{-0.4cm}
 \divider
% \smallskip

 \cvevent{Research Assistant}{Ayers Research Group, McMaster University}{May 2015 -- Aug 2015}{Hamilton, Canada}
% \begin{itemize}
%\item Constructed equations constraining a two-electron reduced density matrix (2-RDM) to represent a many-electron quantum system.
% \begin{itemize}
% \item Implemented a semi-definite optimization algorithm constraining two-electron reduced density matrices (2-RDMs) to represent many-electron quantum systems.
% \end{itemize}


% \newpage
\cvsection{Research}
\footnotesize h-index = 7, i10-index = 5, total citations = 171 (Google Scholar on 25/08/2026)
\begin{flushright}
\vspace{-1.5em} 
\footnotesize \textsuperscript{\textdagger} = These authors contributed equally %\vspace{-2.5em}
\end{flushright}
\normalsize
\cvsubsection{Publications \& Conference Proceedings}

\begin{itemize}
\item \fullcite{furlong2025} \href{doi.org/10.1038/s41467-026-75703-4}{[Link]}
\item \fullcite{dumont2023exploiting} \href{https://www.frontiersin.org/journals/neuroscience/articles/10.3389/fnins.2023.1190515/full}{[Link]}
\item \fullcite{dumont2023biologically} \href{https://compneuro.uwaterloo.ca/files/publications/dumont2023brainsci.pdf}{[Link]}
\item \fullcite{voelker2021simulating} \href{https://compneuro.uwaterloo.ca/files/publications/voelker.2021a.pdf}{[Link]}
\item \fullcite{coleman2021optimal} \href{https://link.springer.com/article/10.1007/s10614-021-10179-6}{[Link]}
\end{itemize}

% 
\divider

\cvsubsection{Conference Proceedings \& Posters}

\begin{itemize}
\item \fullcite{dumont2026}
\item \fullcite{simone2026}\href{https://escholarship.org/uc/item/84x3v8n0}{[Link]}
\item \fullcite{furlong2024biologically} \href{https://compneuro.uwaterloo.ca/files/furlong2024.VSA.sampling.ICANN.pdf}{[Link]}
\item \fullcite{furlong2024recurrent} \href{https://ieeexplore.ieee.org/stamp/stamp.jsp?arnumber=10548051}{[Link]}
\item \fullcite{bartlett2023improving} \href{https://compneuro.uwaterloo.ca/files/ICCM_Bartlett2023.pdf}{[Link]}
\item \fullcite{dumont2022} \href{https://escholarship.org/uc/item/3pf7f9b4}{[Link]}
\item \fullcite{bartlett2022} \href{https://cs.uwaterloo.ca/~jorchard/academic/Bartlett_ICCM2022.pdf}{[Link]}
\item \fullcite{dumont2021} \href{https://static1.squarespace.com/static/6102ca347474c263c40150cd/t/610870a432c9d257a80cace4/1627943077187/Cosyne2021_program_book.pdf}{[Link]}
    \item \fullcite{dumont2020} \href{http://compneuro.uwaterloo.ca/files/publications/dumont.2020.pdf}{[Link]}
\end{itemize}

\divider

\cvsubsection{Manuscripts}
\begin{itemize} 
\item \fullcite{sainz2026neural} \href{https://www.biorxiv.org/content/10.64898/2026.02.18.706604v1}{[Link]}
\end{itemize}

% \divider
 % \pagebreak

%\cvsubsection{Thesis}
%\begin{itemize} 
%\item \fullcite{dumont2025} \href{https://hdl.handle.net/10012/21501}{[Link]}
%\end{itemize}


% \divider

\pagebreak


\cvsubsection{Patents}
\begin{itemize} 
\item \fullcite{voelker2022methods} \href{https://patents.google.com/patent/US20220138382A1/}{[Link]}
\end{itemize}

% \divider 

\cvsection{Teaching and Supervision}
\begin{itemize}
\item \textbf{Masters supervision (2026-present):} 
Supervising a masters student thesis project on learning compositional neural representation for planning and reasoning with tree diffusion. 
\item \textbf{Student supervision (2020):} Supervised an undergraduate research assistant during a four-month project developing a novel algorithm for learning successor features for reinforcement learning.
\item \textbf{Introduction to Computational Mathematics (2022):} Sole instructor for two sections of this course, teaching over 120 students. Planned course content and developed lectures, assignments, and exams. %In the Student Course Perceptions survey, over 80\% of students agreed that I helped them understand the course concepts and fostered a supportive learning environment. %is a rigorous introduction to numerical methods.
    \item \textbf{Logic and Computation (2020):} Served as an Instructional Apprentice for an introductory course in mathematical logic. Led tutorials, graded assignments, and provided assistance during office hours.
    \item \textbf{Teaching Assistantships (2017-2022):} Served as a TA in six undergraduate courses in mathematics and computer science.
\end{itemize}

% \divider

% \cvsection{Skills}
%  \textbf{Programming languages:} Python, MATLAB, R, C++
 
% \textbf{Neural \& Machine Learning Frameworks:} PyTorch, Nengo, Nengo-Loihi, SNNTorch

% \textbf{Scientific Communication \& Leadership:} Workshop organization, peer review, international collaboration

\setlength{\columnsep}{0.07cm} %%%
\cvsection{Academic Service}
% \vspace{-10pt}
\begin{itemize}
    \item \textbf{Leadership \& Organization}
\begin{itemize}
\item Topic Area Leader, Spatiotemporal Dynamics in Neural Computation, Telluride Neuromorphic
AI Workshop (2026)
    \item Topic Area Leader, Language and Thought, Telluride Neuromorphic Cognition Engineering Workshop (2025): Successfully proposed and led a topic area; delivered tutorials on cognitive mapping models; mentored PhD-level researchers through the full project lifecycle from conceptualization to demonstrations for the public.
\item Invited Speaker \& Co-Organizer, Telluride Neuromorphic Cognition Engineering Workshop (2024)
\end{itemize}
\item \textbf{Instruction \& Training}
\begin{itemize}
    \item Instructor, Nengo Summer School (2023, 2024, 2025): Led tutorials and mentored participants developing neural models using the Nengo framework.
\end{itemize}
\item \textbf{Peer Review:} Reviewer for \emph{Nature Communications}, \emph{Communications Biology}, \emph{CogSci 2025}, and \emph{Frontiers in Computational Neuroscience} 
\end{itemize}

% \noindent
% \begin{minipage}{0.55\linewidth}
%     \begin{itemize}
% \item  Co-Organizer of the Language and Thought Topic Area, Telluride Neuromorphic Cognition Engineering Workshop (\textbf{2025})
% \item  Reviewer, Nature Communications (\textbf{2025}) 
% \item  Reviewer, CogSci 2025 (\textbf{2025})

% \end{itemize}
% \end{minipage}
% \begin{minipage}{0.45\linewidth}
% \hfill
% \begin{itemize}
% \item  Invited speaker, Telluride Neuromorphic Cognition Engineering Workshop (\textbf{2024})
% \item  Instructor, Nengo Summer School (\textbf{2024, 2023})
% \item  Reviewer, Frontiers in Neuroscience (\textbf{2024}) %Peer-reviewed two manuscripts as a
% \end{itemize}
% \end{minipage}
% \vspace{0.2em}


% \divider

\cvsection{Awards \& Distinctions}
% \vspace{-10pt}
% \vspace{-1em}
\noindent
\begin{itemize}
\item UZH Postdoc Grant (2026)
\item Marie Skłodowska-Curie Postdoctoral Fellowship (MSCA-PF) Seal of Excellence (2026)
    \item Best New Neuromorph, Telluride Neuromorphic Cognition Engineering Workshop  (2024)
\item European Neural Network Society Best Paper Award (2024)
\item Barbara Hayes-Roth Award for Women in Math and Computer Science (2022)
\item GO-Bell Scholarship (2019–2021)
\item Provost Doctoral Entrance Award for Women (2019)
\item Keith \& Debbie Geddes Graduate Scholarship (2019)
\item University of Waterloo Graduate Scholarship (2018–2019)
\item Emanuel Williams Scholarship in Physics (2014)
\end{itemize}
% \begin{minipage}{0.54\linewidth}
%     \begin{itemize}
%     \item  Best New Neuromorph, Telluride Neuromorphic Cognition Engineering Workshop (\textbf{2024})
% \item  European Neural Network Society  (ENNS) Best Paper Award (\textbf{2024})
%     \item Barbara Hayes-Roth Award for Women in Math and Computer Science (\textbf{2022}) 
%     \end{itemize}
% \end{minipage}
% \begin{minipage}{0.46\linewidth}
% \hfill
% \begin{itemize}
% \item  GO-Bell Scholarship (\textbf{2019-2020}) 
%     \item  Provost Doctoral Entrance Award for Women (\textbf{2019})
%     \item  Keith \& Debbie Geddes Graduate Scholarship (\textbf{2019})
%     \item University of Waterloo Graduate Scholarship (2018, 2019)
%     \item  The Emanuel Williams Scholarship in Physics (\textbf{2014})
% \end{itemize}
% \end{minipage}
    % \end{multicols}


% \noindent
% \begin{minipage}{0.4\linewidth}
%     \begin{itemize}
% \item \textbf{2024} -- European Neural Network Society Best Paper Award  
%     \item \textbf{2022} -- Barbara Hayes-Roth Award for Women in Math and Computer Science 
% \end{itemize}
% \end{minipage}
% \begin{minipage}{0.5\linewidth}
% \hfill
% \begin{itemize}
%  \item \textbf{2020}, \textbf{2019} -- Go-Bell Scholarship 
%     \item \textbf{2019} -- Provost Doctoral Entrance Award for Women 
%     \item \textbf{2019} -- Keith \& Debbie Geddes Graduate Scholarship 
%     %\item University of Waterloo Graduate Scholarship (2018, 2019)
%     \item \textbf{2014} -- The Emanuel Williams Scholarship in Physics
% \end{itemize}
% \end{minipage}





% \cvsection{Academic Projects}
% % \cvproject{Spiking neural network model of simultaneous localization and mapping with Spatial Semantic Pointers}{Poster presentation; Cosyne 2021}
% % \smallskip
% % \begin{itemize}
% % \item Proposed a biological-plausible SLAM model called SSP-SLAM that uses a hybrid oscillatory interference/ continuous attractor network of grid cells for path integration.
% % \end{itemize}

% % \divider


% \cvproject{Successor Features for Transfer Learning in an Advantage Actor-Critic Framework}{Course Project; CS 885: Reinforcement Learning}
% \smallskip
% \begin{itemize}
% \item Implemented an A2C model for reinforcement learning in pytorch that uses a new method for learning successor features called Generalized Successor Estimation.
% \end{itemize}

% \divider

% %\cvproject{Spatial Memory with Semantic Pointers}{SYDE 750: Simulating Neurobiological Systems}
% %\smallskip
% %\begin{itemize}
% %\item Implemented a method for encoding continuous variables with spiking neurons using the Semantic Pointer Architecture and the Nengo library. 
% %\end{itemize}



% \cvproject{Constructing Textual Artificial Conversational Entities using Deep Learning}{Course Project; STAT 841: Statistical Learning - Classification}
% \smallskip
% \begin{itemize}
% \item Built a chatbot using a sequence-to-sequence model with long short-term memory (LSTM) units and an attention mechanism. 
% \end{itemize}
% \smallskip

% \divider

% \cvproject{Dirichlet Mixture Model Library}{Course Project; STAT 840: Computational Inference}
% \smallskip
% \begin{itemize}
% \item Developed a library in Julia (and an R package wrapper) using Dirichlet process mixture models to perform unsupervised, non-parametric clustering.
% \end{itemize}
% \smallskip

% \cvsection{Hobbies and Interests}
% In my spare time, I enjoy getting creative with photography and painting, and camping on the Canadian shield. 

%I  volunteering, and organizing work in the student community.

%\cvproject{Spatial Epidemics Dynamics and Synchronization}{Course Project; Math 4MB3: Mathematical Biology}
%\smallskip
%\begin{itemize}
%\item Implemented a stochastic model of infectious disease spread and investigated the conditions required for spatial synchronization of disease outbreaks. 
%\end{itemize}





\end{fullwidth}
\end{document}
